\begin{tikzpicture}[
  >=Latex,
  line cap=round,
  line join=round,
  font=\small,
  inicial/.style={densely dashed, gray!65, line width=1pt},
  deformada/.style={blue!65!black, line width=1.2pt},
  malla/.style={blue!45!black, line width=0.65pt},
  vector/.style={->, violet!85!black, line width=1.25pt}
]
% Vista esquemática de una cáscara esférica en proyección tridimensional.
\node[font=\small\bfseries] at (0,1.72) {(a) Antes};
\path[fill=blue!7, even odd rule]
  (0,0) ellipse (1.28 and 1.05)
  (0,0) ellipse (0.57 and 0.46);
\draw[inicial] (0,0) ellipse (1.28 and 1.05);
\draw[inicial] (0,0) ellipse (0.57 and 0.46);
% Meridianos y paralelos sugieren la geometría tridimensional.
\draw[inicial] (0,0) ellipse (1.28 and 0.34);
\draw[inicial] (0,0) ellipse (0.48 and 1.02);
\draw[inicial] (0,0) ellipse (0.57 and 0.17);
\draw[inicial] (0,0) ellipse (0.22 and 0.45);
\node[gray!70!black] at (0,-1.43) {cáscara hueca inicial};

\begin{scope}[xshift=6.2cm]
  \node[font=\small\bfseries] at (0,1.72) {(b) Después};
  % El contorno discontinuo conserva la referencia inicial para comparar.
  \draw[inicial] (0,0) ellipse (1.28 and 1.05);
  \draw[inicial] (0,0) ellipse (0.57 and 0.46);
  % Caso ilustrativo: presión interior p_b=p y presión exterior nula.
  \path[fill=blue!10, even odd rule]
    (0,0) ellipse (1.53 and 1.25)
    (0,0) ellipse (0.72 and 0.58);
  \draw[deformada] (0,0) ellipse (1.53 and 1.25);
  \draw[malla] (0,0) ellipse (1.53 and 0.40);
  \draw[malla] (0,0) ellipse (0.58 and 1.22);
  \draw[deformada] (0,0) ellipse (0.72 and 0.58);
  \draw[malla] (0,0) ellipse (0.72 and 0.20);
  \draw[malla] (0,0) ellipse (0.29 and 0.57);
  % Desplazamientos radiales sobre las superficies exterior e interior.
  \draw[vector] (1.29,0) -- (1.51,0);
  \draw[vector] (-1.29,0) -- (-1.51,0);
  \draw[vector] (0,1.06) -- (0,1.23);
  \draw[vector] (0,-1.06) -- (0,-1.23);
  \draw[vector] (0.58,0) -- (0.70,0);
  \draw[vector] (-0.58,0) -- (-0.70,0);
  \node[violet!85!black, anchor=west] at (0.84,0.78) {$\vec u$};
  \node[gray!70!black] at (0,-1.54) {$p_b=p$, $p_a=0$; $\vec u$ radial hacia fuera};
\end{scope}
\end{tikzpicture}